
\hypertarget{bandwidth-complexity}{%
\subsubsection{Bandwidth Complexity}\label{bandwidth-complexity}}

\label{sec:bandwidth-complexity}

\subsubsubsection{Full Node}

What data must a full node download?
A full node must be able to \emph{completely validate} a \emph{single chain}.
For a simplex-chain, provided that the PoRs and corresponding headers remain available (which they always do\footnote{
    The necessary PoRs are, at the very least, part of other simplex-chains, so ``always do'' assumes that simplex-chains themselves remain available, excluding planned shutdown.
    Since all blockchain networks \emph{depend} on the availability of their chains, this is a safe assumption.
}), this means it must download: all blocks for that simplex-chain, and any auxiliary data necessary to verify PoRs.
Note that bandwidth requirements differ based on which UT protocol variant is used.

For \UT{\text{+OP}}, the data a full node requires are: each block, the headers of all reflecting chains, and the missing branches for all PoRs.
Network-wide, headers consume \(N_1 \cdot B_f \cdot B_h\) B/s.
For a single chain, PoRs use \(N_1 \cdot B_f \cdot g \cdot \ceil{\log_2 N_1}\) B/s of capacity, where $g$ is the digest size of the hash used for merkle trees (usually 32 bytes).
Let's denote the total bandwidth required for a full node \(\Delta s\).
In the worst case (\UT{\text{+HO}} and \UT{\text{+HOT}}), where both headers and PoRs must be downloaded:
\begin{align}
    \Delta s & = k_1 + (N_1 \cdot B_f \cdot B_h) + (N_1 \cdot B_f \cdot g \cdot \ceil{\log_2 N_1}) \nonumber \\
    & = k_1 + N_1 \cdot B_f \cdot ( B_h + g \cdot \ceil{\log_2 N_1}) \nonumber
\end{align}

Thus, with \emph{omitted proofs}, \(O(\Delta s) = O(k_1 + N_1 \cdot \log_2 N_1) = O(c \cdot \log_2 c)\).

However, with \emph{explicit PoRs} (variants including +PoRs), \(\Delta s \le k_1 + N_1 \cdot B_f \cdot B_h = O(k_1) = O(c)\).

\aside{
  The term $g \cdot \ceil{\log_2 N_1}$ in these equations is the size of a merkle branch for a PoR.
  What if verkle trees are used instead?
  With a branching factor of 256, 32 byte commitments and proofs, and 1 byte location specifiers, this term should be replaced with $(1+32) \cdot \max(1, \log_{256} N_1)$.
  The complexity of these two terms is the same --- $O(\log c)$ --- but in practice verkle PoRs are less than half the size of merkle PoRs.
  For this reason, the numerical calculations used in the tables throughout this paper assume that the UT implementation uses verkle trees.
}

\subsubsubsection{PoR Graph}

We found the bandwidth requirements for fully reconstructing the PoR graph before in \autoref{sec:por-graph-max-refl} involving the measure of propagation delay, $\phi$ seconds.
Let's call the total bandwidth requirements for the PoR graph $\Delta r$ and substitute in \autoref{eq:simplex-N1}.
\begin{align}{}
    \Delta r &= N_1 B_f (B_h + g(1 + N_1 B_f \phi)) \nonumber \\
    &= \frac{k_1}{2 B_h} \Bigl(B_h + g\Bigl(1 + \frac{k_1 \phi}{2 B_h}\Bigr)\Bigr) \nonumber \\
    % &= \frac{k_1}{2} + \frac{g k_1}{2 B_h}\Bigl(1 + \frac{k_1 \phi}{2 B_h}\Bigr) \nonumber \\
    % &= \frac{k_1}{2} + \frac{g k_1}{2 B_h} + \frac{g {k_1}^2 \phi}{4 {B_h}^2} \nonumber \\
    &= \frac{k_1}{2}\Bigl(1 + \frac{g}{B_h} + \frac{g {k_1} \phi}{2 {B_h}^2}\Bigr) \label{eq:por-graph-delta-r} \\
    \implies O(\Delta r) &= O(c^2\phi) \nonumber
\end{align}
As discussed in \autoref{sec:por-graph-max-refl} it is unclear whether $O(c^2\phi) = O(c)$.
However, we don't need to know $O(\Delta r)$ to estimate some real world values of $\Delta r$.
Practically, $\Delta r$ looks to be between $k_1$ and $2 k_1$ for $k_1=3000$ --- but, this is sensitive to $\phi$.

If we try to restrict $\Delta r < k_1$, then we find that we require a large $B_h$ or small $k_1$:
\begin{align*}
    k_1 &> \frac{k_1}{2}\Bigl(1 + \frac{g}{B_h} + \frac{g {k_1} \phi}{2 {B_h}^2}\Bigr) \\
    2 {B_h}^2 &> 2 B_h g + g k_1 \phi \\
    k_1 &< 2 B_h (B_h - g) (g \phi)^{-1}
\end{align*}
\pz{$\phi$ also used here}
If $k_1 = 3000$, $\phi = 0.5$, we have $B_h \ge 172$ --- much larger than what we've assumed to be the minimum so far.
Achieving $\Delta r < 3000$ requires $\phi \lesssim 0.2$ given the usual parameters.

Regarding the terms from \autoref{eq:por-graph-delta-r} within the parentheses, it is clear that $1 + g/B_h$ will be a value around 1.3, give or take.
What do we expect the value of $g k_1 \phi / (2 {B_h}^2)$ to be?
If $B_h \in [50, 500]$, and $\phi \in [0.05, 2.0]$, then we can see that the coefficient of $k_1$ is, at most, 0.0128 (and $3.2\cdot10^{-7}$ at least).
Thus, this last term dominates $(1+g/B_h)$ when $g k_1 \phi \gtrsim 2 {B_h}^2$.

To keep $\Delta r$ practicable, we should therefore go to some effort to keep $g k_1 \phi \lesssim 2 {B_h}^2$.

\subsubsubsubsubsection{\AxiomOfUnifiedAncestry Optimization}

Using the \AxiomOfUnifiedAncestry, we do not need to reference a block's parents directly, since they are implied by the PoR graph and the absence of fraud proofs.
This results in a small optimization where the parent hash is replaced with the best tip of the longest PoR chain.
Regarding $\Delta r$, $g(1 + N_1 B_f \phi) \to g N_1 B_f \phi$ and thus:
\begin{align}{}
    \Delta r &= N_1 B_f (B_h + g N_1 B_f \phi) \nonumber \\
    % &= \frac{k_1}{2 B_h} \Bigl(B_h + \frac{g k_1 \phi}{2 B_h}\Bigr) \nonumber \\
    &= \frac{k_1}{2}\Bigl(1 + \frac{g {k_1} \phi}{2 {B_h}^2}\Bigr) \nonumber % \label{eq:por-graph-delta-r} \\
\end{align}
Restricting $\Delta r < k_1$ now implies $k_1 < 2 {B_h}^2 (g \phi)^{-1}$.
If $k_1 = 3000$, $\phi = 0.5$, we have $B_h \ge 155$.
So for a maximal simplex of reasonable parameters, we still have $k_1 < \Delta r < 2 k_1$.


\subsubsubsection{Complete Simplex}

\aside{
    This is also the bandwidth required by all active miners due to the \AxiomOfAvailability.
    Miners do not verify the state transitions of each block, but do verify the integrity of the data structure and PoR graph (which is very cheap by comparison).
}

What about the bandwidth required to verify \emph{the entire simplex}?

If miners temporarily keep the blocks of every simplex-chain (so that they can regenerate PoRs and verify that reflected headers correspond to existent blocks) then what is the complexity and burden of this?
% Each simplex-chain has a raw throughput of \(k_1\) bytes/s.
% From \autoref{eq:simplex-N1} we know that \(N_1 = \frac{k_1}{2 \cdot B_f \cdot B_h}\).

The amount of network bandwidth, \(\Delta S\), required to download all blocks (as they are produced) across all simplex-chains is the overall transaction throughput, $(N_1 \cdot k_{1,tx})$, plus the data necessary to reconstruct the PoR graph: $\Delta r$.
Any auxiliary data can be deterministically regenerated.
\begin{equation}
\begin{split}
\Delta S &= N_1 \cdot k_{1,tx} + \Delta r
\label{eq:bandwidth-req} \\
% \\ & = \frac{{k_1}^2}{2 \cdot B_f \cdot B_h}
\implies O(\Delta S) &= O(c^2) \nonumber
\end{split}
\end{equation}

It is clear that \(\Delta S\) has order \(O(c^2)\), but how bad is this?
For \(k_1 = 3000\), \(B_f = \frac{1}{60}\), and \(B_h = 112\): \(\Delta S \approx 1.2\) MB/s.
With those figures: \(N_1 \approx 800\) simplex-chains, \(N_2 \approx 645,000\) dapp-chains, and maximum tps of \(\sim 7.7\times 10^{6}\) at nesting level 2.
Decreasing block times to 15s correspondingly decreases the bandwidth requirements to 0.3 MB/s for a simplex with \(\sim 200\) chains, \(\sim 40,000\) dapp-chains, and \(\sim 484,000\) max tps.

While \(O(c^2)\) bandwidth scaling is not ideal, it's clear that --- especially in the early days of a UT simplex when there are fewer simplex-chains --- there are tolerable configurations available; i.e., there is \emph{excess capacity}.
